% Synthetic placeholder figure generated from experiment/sec-5-rq5-extensibility/data/new_capability_eval.csv
\begin{figure}[t]
  \centering
  \begin{tikzpicture}
    \begin{groupplot}[
      group style={group size=2 by 1, horizontal sep=1.2cm},
      bpfplot,
      width=0.45\linewidth,
      height=0.62\linewidth,
      xbar,
      ytick=data,
      y dir=reverse,
      symbolic y coords={ipv4 parse,Katran parser,Tracee event,Cilium svc key},
      enlarge y limits=0.18
    ]
      \nextgroupplot[
        title={(a) Speedup},
        xmin=1.0,
        xmax=1.20,
        xlabel={Execution speedup over stock}
      ]
        \addplot[fill=bpfBlue!75, draw=bpfBlue!75] coordinates {
          (1.17,ipv4 parse)
          (1.10,Katran parser)
          (1.07,Tracee event)
          (1.04,Cilium svc key)
        };

      \nextgroupplot[
        title={(b) Code shrink},
        xmin=0,
        xmax=12,
        xlabel={Code-size reduction (\%)}
      ]
        \addplot[fill=bpfOrange!80, draw=bpfOrange!80] coordinates {
          (11.2,ipv4 parse)
          (8.6,Katran parser)
          (5.1,Tracee event)
          (3.7,Cilium svc key)
        };
    \end{groupplot}
  \end{tikzpicture}
  \caption{Benefit of one newly added capability. In this synthetic placeholder example, the new endian-fusion family yields a clear win on its dedicated microbenchmark and smaller but still visible gains on real programs. \note{rq-5-exp-2}}
  \label{fig:rq5-new-capability}
\end{figure}
